\section{Introduction}\label{sec:introduction}
Let $V=\{1,\ldots,n\}$, where $n\ge4$, and let $d=(d_1,\ldots,d_n)$ be a
graphical degree vector. Write $\Omega_d$ for the set of simple undirected
graphs on $V$ whose labeled degree vector is $d$, put $M=|\Omega_d|$, and let
$\pi$ be the uniform probability measure on $\Omega_d$. A \emph{switch}
removes two disjoint edges and inserts a different perfect matching on their
four vertices, provided both new edges are absent. A \emph{Curveball trade} at
a vertex pair $\{i,j\}$ redistributes the vertices adjacent to exactly one of
$i$ and $j$ between these two endpoints, keeping both endpoint degrees and
every other adjacency. Both moves preserve simplicity and every vertex degree.

We use the following precise discrete-time conventions. The switch chain $P$
holds with probability $1/2$. Otherwise it chooses a four-set uniformly and one
of the six ordered pairs of distinct perfect matchings on that set uniformly.
It removes the first matching and adds the second if all removed edges are
present and all added edges absent; otherwise it holds. The Curveball chain
$K$ chooses an unordered vertex pair uniformly and moves to a graph chosen
uniformly among those obtainable from the current one by a Curveball trade at
that pair, the current graph included; no additional holding probability is
imposed. Time
includes rejected proposals and holds for $P$, and counts trades for $K$. For
probability measures on $\Omega_d$ and $Q\in\{P,K\}$, set
\[
 \|\mu-\nu\|_{\mathrm{TV}}
 =\frac12\sum_{G\in\Omega_d}|\mu(G)-\nu(G)|,
 \qquad
 t_{\mathrm{mix}}^Q(\varepsilon)
 =\min\Bigl\{t\in\mathbb Z_{\ge0}:
 \max_{G\in\Omega_d}\|Q^t(G,\cdot)-\pi\|_{\mathrm{TV}}\le\varepsilon\Bigr\}.
\]
For $M>1$ write $\operatorname{gap}(Q)=1-\lambda_2(Q)$, where $\lambda_2(Q)$ is
the second largest eigenvalue of $Q$, which is symmetric
(Section~\ref{sec:setup}). Here and throughout, logarithms are
natural.

\begin{theorem}\label{thm:main}
For every graphical $d$ on $n\ge4$ vertices with $M>1$,
\begin{equation}\label{eq:gaps}
 \operatorname{gap}(P)\ge\frac{1}{6\binom n4},\qquad
 \operatorname{gap}(K)\ge\frac{2}{\binom n2}.
\end{equation}
For every $n\ge4$, both constants are best possible universal constants,
that is, the worst case over $d$: equality holds in both bounds for
$d=(2,2,1,1,0,\ldots,0)$.
\end{theorem}

Equality holds whenever $M=2$. Section~\ref{shp:section} proves this,
describes the extremal instances, and states what is not claimed for larger
realization spaces.

\begin{corollary}[Mixing times]\label{cor:mixing}
For $M>1$ and $0<\varepsilon\le1/2$, set
$A_\varepsilon=\log(1/(2\varepsilon))+\tfrac12\log(M-1)$. Then
\begin{align}
 t_{\mathrm{mix}}^P(\varepsilon)&\le\left\lceil6\binom n4 A_\varepsilon\right\rceil,
 &t_{\mathrm{mix}}^K(\varepsilon)&\le\left\lceil\frac{\binom n2}{2}
 A_\varepsilon\right\rceil,\label{eq:mixing}\\
 t_{\mathrm{mix}}^P(1/4)&\le n^6/4+1,
 &t_{\mathrm{mix}}^K(1/4)&\le n^4/16+1.\label{eq:coarse}
\end{align}
If $M=1$, both mixing times are zero.
\end{corollary}

No sharpness of these mixing bounds is asserted.

\paragraph{History and significance.}
Uniform sampling of graphs with prescribed degrees is a basic problem for
degree-constrained network models. Kannan, Tetali and Vempala~\cite{KTV99}
proved rapid mixing of the switch chain for near-regular bipartite degree
sequences and proposed the general rapid-mixing problem. Cooper, Dyer and
Greenhill~\cite{CDG07,CDG12} established polynomial mixing for regular
undirected graphs. Greenhill and Sfragara~\cite{GS18} extended this to
irregular undirected degrees under sparsity and degree conditions, and treated
directed degrees under additional hypotheses including irreducibility.
Erd\H{o}s, Greenhill, Mezei, Mikl\'os, Solt\'esz and Soukup~\cite{Unified22}
proved rapid mixing on $P$-stable families of undirected, bipartite and
directed degree sequences. These results impose restrictions on the degree
families. Tikhomirov and Youssef~\cite{TY23} proved sharp Poincar\'e bounds
for regular bipartite graphs in a specified degree range and sharp
log-Sobolev bounds for fixed degree; their setting is distinct from arbitrary
undirected degrees.

Verhelst~\cite{Verhelst08} developed trade-based MCMC sampling for binary
matrices with fixed marginals, including a uniform-stationary variant.
Carstens, Berger and Strona~\cite{CBS18} defined undirected Curveball trades
and proved convergence to the uniform distribution. Their mixing-rate
comparisons are experimental, and they left theoretical rapid mixing open.
Carstens and Kleer~\cite{CK18} used the Johnson spectrum to compare switch and
Curveball relaxation times for binary matrices, including irreducible spaces
with forbidden entries. Their comparison transfers rapid mixing from the
switch chain; it supplies no universal absolute gap, and the undirected
comparison is left open in their preprint.

The closest universal results are those of OpenAI~\cite{OAI26} and of Fu, Qin
and Wang~\cite{FQW26}. The manuscript~\cite{OAI26} is machine-checked in Lean
and has not been peer reviewed. It develops the pair-fiber framework used
here, with the triple reduction and the disjoint-cell and slice Gram methods,
and proves that $t_{\mathrm{mix}}^P(1/4)\le2n^8$ and that $P$ has spectral gap
at least $[24n^2\binom n4]^{-1}$ \cite[Theorem~1.1]{OAI26}. Fu, Qin and Wang prove sharp swap-chain and
row-pair resampling bounds for all feasible binary matrix margins; their
Theorem~5.1 also gives the three-row input used here. Arbitrary row-pair moves
do not preserve symmetric zero-diagonal adjacency matrices, so their global
matrix bound does not directly imply Theorem~\ref{thm:main}.
Theorem~\ref{thm:main} replaces the gap bound of \cite[Theorem~1.1]{OAI26} by
the best possible universal constant for the same chain, and gives the best
possible universal constant for undirected Curveball.

\paragraph{Proof strategy.}
For an unordered vertex pair $a$, let $E_a$ average a function over the
graphs reachable by a Curveball trade at $a$, put $h_a=I-E_a$ and $H=\sum_ah_a$, and let $L$ be the unnormalized
Laplacian of the switch graph on $\Omega_d$. Then $I-P=L/(12\binom n4)$ and
$K=I-H/\binom n2$ (Section~\ref{sec:setup}), so Theorem~\ref{thm:main} reduces
to $L\succeq H\succeq2(I-\Pi)$, where $\succeq$ is the order of quadratic
forms and $\Pi$ is projection onto constants. The comparison $H\preceq L$ is
the classical Johnson comparison. The second inequality follows from
$H^2\succeq2H$ and $\ker H=\mathbb R\mathbf 1$. As in
\cite[Proposition~7.1]{OAI26}, we expand
\[
 H^2=\sum_{|T|=3}H_T^2-(n-3)H+D,
\]
where $H_T$ sums $h_a$ over the three pairs $a$ inside a vertex triple $T$,
and $D$ is the
ordered sum of the products $h_ah_b$ over disjoint pairs $a,b$. OpenAI
proved $H_T^2\succeq H_T$ and $D\succeq0$, hence $H^2\succeq H$
\cite[Propositions~4.1 and~7.1, Corollary~6.3]{OAI26}. That intermediate
estimate, with the constant kernel and the classical Johnson comparison,
already gives $\operatorname{gap}(P)\ge1/(12\binom n4)$ and
$\operatorname{gap}(K)\ge1/\binom n2$, and, by the argument of
Corollary~\ref{cor:mixing}, $t_{\mathrm{mix}}^P(1/4)\le n^6/2+1$. In
particular the inverse-fourth order follows from the ingredients
of~\cite{OAI26} plus the Johnson comparison; it is not the new part of this
paper, and the improvement of Corollary~\ref{cor:mixing} over those
ingredients is the factor two.

The new step is the weighted estimate $D\succeq H$, which gives
$H^2\succeq(n-2)H-(n-3)H+H=2H$. It yields the exact factor two and the
Curveball constant $2/\binom n2$; Section~\ref{shp:section} shows separately
that both constants are attained. We prove $D\succeq H$ in two weighted
steps. For each disjoint pair, the interaction $h_ah_b+h_bh_a$ is
bounded below by twice a cross-square projection minus projections onto
centered equality indicators, each weighted by the probability $p_a$ that two
singleton labels of its fiber receive the same endpoint. On every slice, the
Johnson energy minus the weighted equality projections still dominates the
centered identity. Without the weights this slice inequality fails
(Remark~\ref{asm:weighting}). The cell analysis follows
\cite[Proposition~5.1]{OAI26} up to its final square, and the slice analysis
reuses the Gram decomposition of \cite[Lemma~6.1]{OAI26}. New here are the
weights $p_a$ with the completed square in the cell estimate, the exact
second slice-moment identity, the weighted slice inequality $G-pS\succeq h$,
and the assembly $D\succeq H$.

\paragraph{Organization.}
Section~\ref{sec:setup} defines the pair fibers and operators and identifies
both chains with them. Section~\ref{joh:section} proves the Johnson gap and
$H\preceq L$. Section~\ref{cell:section} proves the weighted inequality for
one disjoint pair, and Section~\ref{sl:section} the weighted slice
inequality. Section~\ref{asm:section} proves connectivity, $D\succeq H$ and
$H^2\succeq2H$, and tracks the constants through Theorem~\ref{thm:main} and
Corollary~\ref{cor:mixing}. Section~\ref{shp:section} proves equality
whenever $M=2$ and gives extremal instances. Section~\ref{fv:section} describes the Lean verification, and
Section~\ref{op:section} states open problems. Appendix~\ref{app:section}
proves the three-row and graph-triple inputs. All proof ingredients are
supplied in the paper.
